% sections/04_saddle_to_saddle.tex -- The saddle-to-saddle regime (Q1-Q4, H1-H3). SKELETON.
% Main section of the report. Sources: PLAN.md Q1-Q4; briefs s2s-core, order (and locality);
% results in experiments/s2s_core/, experiments/locality/, experiments/order/ and the notes
% notes/s2s_core.md, notes/locality.md, notes/order.md; theory in notes/theory_escape.md and
% notes/theory_saddles.md; pilot experiments/pilot_h2_locality/.
\section{The saddle-to-saddle regime}\label{sec:s2s}

\placeholder{Skeleton. Theory and experiments for gradient flow from small initialization:
plateau sequence, escape times and their dependence on each parameter, locality of later escapes,
learning order, saddle types.}

\subsection{Phenomenology}\label{sec:s2s-phenomenology}
\placeholder{Representative runs: loss in linear and logarithmic time, overlaps, bond spectra on
every edge, effective ranks; S2S signature (\cref{sec:setup-diagnostics}).}
% TODO(Q1): fan-in init, odeco teachers with separated weights [3, 2, 1] and close weights
% [1, 0.85, 0.7]; binary N = 4, 8, MPS N = 4, 8, HT N = 4; chi = R and 2R; 10 seeds.
% Questions: are plateau levels the partial-fit levels L_k (eq:partial-fit)? Do bond spectra grow
% one value at a time on every edge simultaneously? Are the plateaus near critical points of the
% bond-rank-constrained problem? Status expected: \empirical. PLAN.md: "partly known" (Razin et al.
% 2022 for hierarchical tensor factorization; Jacot et al. 2021 for DLNs).

\subsection{An exactly solvable case: aligned initialization}\label{sec:s2s-aligned}
\placeholder{\Cref{prop:aligned} in experiments: TTN gradient flow against the exact mode
solutions for binary, MPS and HT trees.}
% TODO: overlay figure, max deviation. This reproduces Saxe et al. (2014) in the TTN setting:
% mark \knownresult[\citealp{saxe2014exact}]; it is a demonstration and integrator check, not a
% new claim. Include the escape-time asymptotics eq:aligned-escape and the unbalanced reduction
% (ttn/theory.py: m smallest cores give a depth-m mode).

\subsection{Escape from the origin}\label{sec:s2s-escape}
\placeholder{First escape time $t_1$ versus $\isc$, $\chi$ and $\Kc$, against the prediction.}
% TODO(Q2): hypothesis t_1 ~ rho_0^{-(K-2)} [tau_align + K^{K/2-1} / ((K-2) s_1)], rho_0 = ||theta_0||
% (spherical flow of the degree-K homogeneous g = <W, W*> plus radial growth; Kumar & Haupt-type
% directional convergence). Consequences: t_1 ~ sigma^{-(K-2)}; at fixed sigma t_1 decreases with
% chi through rho_0^2 ~ sigma^2 sum_v chi_v (Prop. init-norm). Figures: t_1 vs sigma per tree with
% fits and theory; prefactor distribution over seeds vs theory; t_1 vs chi; t_1 vs K. Table of
% fitted exponents (with s.e.) vs K - 2. Separate parameter-free predictions from fitted ones.
% Status: \derived (theory, under the stated assumptions) + \empirical (fits); \newresult for TTNs.

\subsection{Later saddles and locality}\label{sec:s2s-locality}
\placeholder{Escape times after the first plateau depend on where the new component lives in the
tree.}
% TODO(Q3, H2): k = number of cores in the Steiner subtree of the sites where the new component
% differs from what is learned (Prop. teacher-spectra (b)). Prediction: t_{i+1} - t_i ~
% sigma^{-(k-2)} for k >= 3, ~ log(1/sigma)/lambda for k = 2, no plateau for k = 1. The same
% teacher has different saddle sequences on an MPS and on a binary tree. Pattern teachers (local,
% far, global) on MPS, binary, HT, custom trees, N = 4-8. Exponent table: measured vs predicted k.
% (H3) strict saddle iff k = 2: gauge-projected Hessian at plateau midpoints (number of negative
% eigenvalues, lambda_min/lambda_max). Pilot: experiments/pilot_h2_locality. Status: believed
% open -> \newresult if confirmed; \conjecture until then.

\subsection{Order of learning}\label{sec:s2s-order}
\placeholder{Which component is learned first, how random the order is, and how a wrong order is
corrected.}
% TODO(Q4, H1, H1b): P(larger first) vs s2/s1 in {0.5, ..., 0.99}, K, chi, d with binomial error
% bars (>= 200 seeds); race-model / basin prediction overlay; matrix-like teacher as control
% (greedy). H1b: after the smaller component is learned first, a strict saddle with Hessian
% eigenvalue q (s_learned - s_new) < 0 rotates the learned component in time O(log 1/sigma).
% Three or more components: Kendall tau with the greedy order. Known ingredients: odeco components
% are attracting fixed points of the tensor power method (anandkumar2014tensor); directional
% convergence near small init (Kumar & Haupt).

\subsection{Dependence on each parameter}\label{sec:s2s-parameters}
\placeholder{Summary table: initialization scale, number of cores and tree shape, bond dimension,
physical dimension, teacher spectrum, imbalance, learning rate; effect on plateau levels, durations
and order; status of each statement.}
% TODO: one row per parameter with: effect on the saddle sequence | on plateau durations | on the
% order | evidence (figure) | status tag | new/known.
